% Generated draft: requires mathematical and source review.
\documentclass[11pt]{article}
\usepackage[margin=25mm]{geometry}
\usepackage{fontspec}
\setmainfont{Latin Modern Roman}
\usepackage{amsmath,amssymb,mathtools,unicode-math}
\setmathfont{Latin Modern Math}
\usepackage{enumitem,needspace,xurl,hyperref}
\setlength{\parindent}{0pt}
\setlength{\parskip}{5pt}
\setlength{\emergencystretch}{3em}
\newcommand{\N}{\mathbb{N}}
\newcommand{\Z}{\mathbb{Z}}
\newcommand{\Q}{\mathbb{Q}}
\newcommand{\R}{\mathbb{R}}
\newcommand{\C}{\mathbb{C}}
\newcommand{\F}{\mathbb{F}}
\newcommand{\A}{\mathbb{A}}
\newcommand{\PP}{\mathbb{P}}
\newcommand{\E}{\mathbb{E}}
\newcommand{\eps}{\varepsilon}
\newcommand{\dd}{\,\mathrm{d}}
\begin{document}
\section*{UnsolvedMath ID: 701089}
\textbf{Source problem code:} catalogue entry
\subsection{Definitions and mathematical statement}
Let \(n\geq 1\) be an integer, and let \(w=w(x_1,\ldots,x_n)\) be a group word, that is, an element of the free group on \(x_1,\ldots,x_n\). For a group \(H\), let \(w(h_1,\ldots,h_n)\) denote evaluation of this word at a tuple in \(H^n\). Define
\[
V(w)=\{H: w(h_1,\ldots,h_n)=1_H\text{ for every }(h_1,\ldots,h_n)\in H^n\}.
\]
Thus \(V(w)\) is the variety of groups satisfying the law \(w=1\).

Define \(V(w^*)\) to be the class of groups \(G\) such that, for every choice of infinite subsets \(S_1,\ldots,S_n\subseteq G\), there exist elements \(s_i\in S_i\), for \(1\leq i\leq n\), satisfying
\[
w(s_1,\ldots,s_n)=1_G.
\]
Define \(V(w^{\#})\) to be the class of groups \(G\) such that, for every choice of infinite subsets \(S_1,\ldots,S_n\subseteq G\), there exist elements \(s_i\in S_i\), for \(1\leq i\leq n\), with
\[
\langle s_1,\ldots,s_n\rangle\in V(w).
\]
Here \(\langle s_1,\ldots,s_n\rangle\) is the subgroup generated by the selected elements. Its membership in \(V(w)\) requires \(w(h_1,\ldots,h_n)=1_G\) for every tuple of elements of that subgroup. No distinctness or disjointness condition is imposed on the subsets or the selected elements. The symbols \(w^*\) and \(w^{\#}\) label these classes; they are not additional group words.

\begin{enumerate}
\item[(a)] Does there exist an integer \(n\geq 1\), a group word \(w\) in \(n\) variables, and an infinite group \(G\) such that
\[
G\in V(w^{\#})\setminus V(w)?
\]
\item[(b)] Does there exist an integer \(n\geq 1\), a group word \(w\) in \(n\) variables, and an infinite group \(G\) such that
\[
G\in V(w^*)\setminus V(w^{\#})?
\]
\end{enumerate}
The witnesses in the two questions need not be the same.

The supplied source expresses the expectation that both answers are yes. It also reports that no example in part (a) can use any of several words, including
\[
x_1^n,\qquad [x_1,\ldots,x_n],\qquad [x_1,x_2]^2,
\qquad (x_1x_2)^3x_2^{-3}x_1^{-3},
\qquad x_1^{a_1}\cdots x_n^{a_n},
\]
where \(a_1,\ldots,a_n\) are nonzero integers in the final family. This is a supplied background assertion, not an independently verified result. The list retains the source's notation: the parameter ranges for its power and commutator families and the bracketing convention for the iterated commutator are not explicitly specified in the record. The registry labels the problem open; this label is supplied curation rather than an independently verified assessment.
\subsection{Short English statement}
For each group word, compare satisfying its law everywhere with two weaker requirements on arbitrary infinite subsets: selecting elements whose generated subgroup satisfies the law, and selecting elements on which the word itself evaluates to the identity. The two questions ask whether each successive weakening admits an infinite group excluded by the preceding condition. The source expects affirmative answers and reports several excluded word families for part (a), subject to the qualifications above.
\subsection{Sources}
\begin{itemize}
\item Local UnsolvedMath registry, permanent ID \(701089\), exact source problem code \texttt{catalogue entry}, title \emph{Kourovka Notebook Problem 15.1}. The supplied record labels the problem open and records its update date as \texttt{2026-10-16T14:00:00Z}; this is registry metadata, not independent literature verification.
\item Supplied source attribution: \emph{The Kourovka Notebook (Unsolved Problems in Group Theory)}, original problem 15.1, 15th issue (2002); also identified in the record with the 21st issue edition (2026), arXiv identifier \texttt{1401.0300v48}. The collection description names E. I. Khukhro and V. D. Mazurov as editors. These bibliographic details are supplied and have not been independently checked.
\item Attribution discrepancy: the background and proposer field name A. Abdollahi, while the statement begins with the names P. Longobardi, M. Maj, and A. H. Rhemtulla. Both attributions are retained; their respective roles are unresolved in the supplied record.
\item The record supplies no separate references for its reported exclusions in part (a), and no literature-verification date is supplied.
\end{itemize}
\end{document}
